\documentclass[12pt]{article}
\usepackage[T1]{fontenc}
\usepackage[utf8]{inputenc}
\usepackage{lmodern}
\usepackage{textcomp}
\usepackage{enumitem}
\usepackage{geometry}
\geometry{margin=1in}
\begin{document}
\begin{center}
Answer Key - Mathematics - Undergraduate Level Assessment 22 \
\small (Answers are ordered exactly as in the question paper.)
\end{center}

\section*{Multiple Choice}
\begin{enumerate}[label=\arabic*.]
\item \textbf{Answer:} B
\\ \textit{Options:} A) True, True; B) False, False; C) True, False; D) False, True
\\ \textit{Note:} The external direct product of cyclic groups is not always cyclic, as it depends on the orders of the groups involved, and the external direct product of D\_3 and D\_4 is not isomorphic to D\_12 due to their differing structures and orders.
\item \textbf{Answer:} C
\\ \textit{Options:} A) True, True; B) False, False; C) True, False; D) False, True
\\ \textit{Note:} Statement 1 is true because the polynomial ring R[x] inherits the integral domain property from R, while Statement 2 is false because the degree of the product of two polynomials is equal to the sum of their degrees only when neither polynomial is the zero polynomial.
\item \textbf{Answer:} D
\\ \textit{Options:} A) (S, +, x) is closed under addition modulo 10.; B) (S, +, x) is closed under multiplication modulo 10.; C) (S, +, x) has an identity under addition modulo 10.; D) (S, +, x) has no identity under multiplication modulo 10.
\\ \textit{Note:} The statement is correct because the subset S = \{0, 2, 4, 6, 8\} does not contain an element that acts as a multiplicative identity (1) under multiplication modulo 10, thus confirming that (S, +, x) lacks an identity with respect to multiplication.
\item \textbf{Answer:} C
\\ \textit{Options:} A) 27; B) 28; C) 35; D) 37
\\ \textit{Note:} The order of the group G must be a multiple of the order of the subgroup (7) and must satisfy the condition that the group has elements of order 5 (as 35 is the product of 7 and 5, and no element can be its own inverse implies that all non-identity elements must have orders greater than 2).
\item \textbf{Answer:} C
\\ \textit{Options:} A) 4; B) 6; C) 12; D) 24
\\ \textit{Note:} The maximum possible order of an element in the direct product of two cyclic groups, such as Z\_4 and Z\_6, is the least common multiple of their orders, which is lcm(4, 6) = 12.
\end{enumerate}

\section*{True / False}
\begin{enumerate}[label=\arabic*.]
\setcounter{enumi}{5}
\item \textbf{Answer:} False
\\ \textit{Options:} A) True; B) False
\\ \textit{Note:} The statement is false because there exist finitely generated torsion-free abelian groups, such as the group of integers \( \mathbb{Z}^n \) for \( n > 1 \), which are indeed free abelian groups.
\item \textbf{Answer:} False
\\ \textit{Options:} A) True; B) False
\\ \textit{Note:} The statement is false because a subset H of a group G does not necessarily contain the products of its elements with all elements of G, as subsets do not have to be closed under the group operation.
\item \textbf{Answer:} False
\\ \textit{Options:} A) True; B) False
\\ \textit{Note:} The statement is false because a relation is symmetric only if for every pair (a, b) in the relation, the pair (b, a) is also included, which is not necessarily true for all relations.
\item \textbf{Answer:} False
\\ \textit{Options:} A) True; B) False
\\ \textit{Note:} The statement is false because a linear transformation can be expressed as f(x, y) = ax + by, and given the values f(1, 1) = 1 and f(-1, 0) = 2, it cannot be determined that f(3, 5) equals 8 without additional information about the coefficients a and b.
\item \textbf{Answer:} False
\\ \textit{Options:} A) True; B) False
\\ \textit{Note:} The statement is false because if C is false, it does not imply any conditions on A and B, meaning it is possible for both A and B to be true, which contradicts the assertion that this scenario is impossible.
\end{enumerate}

\section*{Long Form Response}
\begin{enumerate}[label=\arabic*.]
\setcounter{enumi}{10}
\item \textbf{Answer:} Let \[t = \frac{x - 2}{3} = \frac{y + 1}{4} = \frac{z - 2}{12}.\]Then $x = 3 t + 2,$ $y = 4 t - 1,$ and $z = 12 t + 2.$ Substituting into $x - y + z = 5$ gives us \[(3 t + 2) - (4 t - 1) + (12 t + 2) = 5.\]Solving, we find $t = 0.$ Thus, $(x,y,z) = \boxed{(2,-1,2)}.$
\\ \textit{Note:} correct because it accurately uses a parameterization of the line to express \(x\), \(y\), and \(z\) in terms of \(t\), substitutes these expressions into the plane equation to solve for \(t\), and then finds the corresponding point of intersection by substituting \(t\) back into the parameterization.
\item \textbf{Answer:} First, we rearrange the equation so that one side is a quadratic written in the usual way and the other side is $0$. We can do this by subtracting $20 x$ from each side (and reordering the terms): $5 x^2-20 x+18 = 0$$This doesn't factor in any obvious way, so we apply the quadratic formula, which gives \begin{align*} x = \frac{-(-20)\pm \sqrt{(-20)^2-4(5)(18)}}{2(5)} &= \frac{20\pm \sqrt{400-360}}{10} \\ &= \frac{20\pm \sqrt{40}}{10} \\ &= 2\pm \frac{\sqrt{40}}{10}. \end{align*}We observe that $\ensuremath{\sqrt{40}}$ is between $6$ and $7$, so $\ensuremath{\frac}\{\ensuremath{\sqrt{40}}\}\{10\}$ is between $0.6$ and $0.7$. Thus one of our solutions is between $1.3$ and $1.4$, while the other is between $2.6$ and $2.7$. Rounding each solution to the nearest integer, we get $1$ and $3$, the product of which is $\ensuremath{\boxed{3}}\$.
\\ \textit{Note:} The answer correctly identifies and rearranges the quadratic equation, applies the quadratic formula to find the distinct solutions for \( x \), and sets the stage for rounding these solutions to the nearest integers before multiplying them, which is the required process to solve the problem.
\end{enumerate}

\vfill
\noindent\textit{End of Answer Key}
\end{document}